\begin{table}[H]
\centering
\caption{Ranking departamental de niveles de productividad laboral, 2024pr}
\label{tab:pib_geih_productividad_departamento_ranking_niveles}
\scriptsize
\begin{tabular}{lrrrr}
\toprule
Departamento & PIB/trab. 2024pr & Puesto & PIB/hora 2024pr & Puesto \\
\midrule
Bogotá D.C. & 67,4 & 1 & 28,6 & 1 \\
Santander & 62,0 & 3 & 26,1 & 2 \\
Meta & 63,1 & 2 & 25,9 & 3 \\
Valle del Cauca & 54,5 & 4 & 23,4 & 4 \\
Boyacá & 47,5 & 6 & 22,6 & 5 \\
Antioquia & 47,7 & 5 & 20,3 & 6 \\
Bolívar & 45,1 & 7 & 19,8 & 7 \\
Cundinamarca & 40,4 & 8 & 17,4 & 8 \\
Tolima & 40,0 & 9 & 16,9 & 9 \\
Risaralda & 38,7 & 10 & 16,4 & 10 \\
Atlántico & 37,3 & 11 & 16,0 & 11 \\
Caldas & 36,6 & 12 & 15,6 & 12 \\
Quindío & 33,8 & 13 & 15,1 & 13 \\
Huila & 31,6 & 15 & 14,2 & 14 \\
Cesar & 32,6 & 14 & 13,8 & 15 \\
Cauca & 26,6 & 17 & 13,7 & 16 \\
Chocó & 28,6 & 16 & 11,9 & 17 \\
La Guajira & 24,8 & 19 & 11,1 & 18 \\
Córdoba & 22,0 & 22 & 10,7 & 19 \\
Caquetá & 24,9 & 18 & 10,5 & 20 \\
Sucre & 22,0 & 23 & 10,4 & 21 \\
Magdalena & 23,9 & 21 & 10,3 & 22 \\
Norte de Santander & 24,3 & 20 & 10,0 & 23 \\
Nariño & 17,8 & 24 & 9,1 & 24 \\
\bottomrule
\end{tabular}
\caption*{\footnotesize Nota: PIB por trabajador en millones de pesos constantes de 2015; PIB por hora en miles de pesos constantes de 2015. La tabla está ordenada por el nivel de PIB por hora trabajada en 2024pr. Fuente: cálculos propios con DANE y GEIH.}
\end{table}